\section{Kimi K2.5: fusión de tres dimensiones}

\subsection{Entrenamiento y arquitectura}

K2.5 fusiona los avances de las tres dimensiones anteriores en un solo modelo:
\begin{itemize}
    \item \textbf{Optimizador Muon + QK-Clip}: mejora la token efficiency
    \item \textbf{Kimi Delta Attention + arquitectura Linear}: refuerza la capacidad de long context
    \item \textbf{Agent Swarms}: abre una nueva dimensión de parallel scaling
\end{itemize}

El modelo se entrenó en GPU NVIDIA H800, con cada nodo con 2 TB de RAM y las GPU interconectadas mediante NVLink. El modelo base de K2.5 se entrenó con más de 30 billones de tokens (base model 15T + 15T adicionales de K2.5), y el proceso de entrenamiento fue extremadamente estable: no hubo ningún loss spike.

\begin{importantbox}{Early Fusion: entrenamiento conjunto nativo de visión y texto}
K2.5 es el primer modelo abierto con \textbf{capacidad nativa conjunta de visión y texto}. A diferencia de los esquemas de late fusion, que agregan capacidad visual sobre un modelo de texto ya existente, K2.5 fusiona los tokens de visión y de texto desde el día 0 del entrenamiento (early fusion); los experimentos iniciales muestran que este esquema supera al de late fusion.
\end{importantbox}

\subsection{Efecto de refuerzo entre modalidades}

Durante el entrenamiento de K2.5 se descubrió un fenómeno interesante: las dos modalidades pueden reforzarse mutuamente.

\begin{importantbox}{Vision mejora Text, Text mejora Vision}
\begin{itemize}
    \item \textbf{Vision $\rightarrow$ Text}: usar RL únicamente con tareas visuales (sin ninguna tarea de matemáticas o programación) logra, aun así, mejorar el desempeño en razonamiento textual
    \item \textbf{Text $\rightarrow$ Vision}: con una base textual sólida, no se necesitan en absoluto datos de SFT visual (Zero Vision SFT); basta con SFT textual + RL conjunto para alcanzar un desempeño visual cercano al SOTA
\end{itemize}
Este hallazgo muestra que el early fusion hace que las dos modalidades compartan de verdad el espacio de representación, lo que permite una transferencia de capacidades en ambos sentidos.
\end{importantbox}

